% =====================================================
% Merged header (Bootcamp palette + Metropolis look)
% =====================================================
\documentclass[10pt,aspectratio=169]{beamer}

% -----------------------------
% Language & encoding
% -----------------------------
\usepackage[utf8]{inputenc}
\usepackage[T1]{fontenc}
\usepackage[english,french]{babel}
\usepackage{lmodern}
\usepackage{microtype}
\usepackage{xcolor}

% -----------------------------
% Math / tables / layout
% -----------------------------
\usepackage{amsmath,amssymb,bm,mathtools}
\usepackage{booktabs}
\usepackage{multicol}
\usepackage{changepage}

% -----------------------------
% Graphics
% -----------------------------
\usepackage{graphicx}
\usepackage{tikz}
\usetikzlibrary{arrows.meta,positioning,calc,decorations.pathreplacing}

\usepackage{pgfplots}
\pgfplotsset{compat=1.18}

% -----------------------------
% Theme (Metropolis)
% -----------------------------
\usetheme[numbering=fraction,progressbar=frametitle]{metropolis}
\setbeamertemplate{frame numbering}[fraction]
\setbeamertemplate{navigation symbols}{}

% -----------------------------
% Colors (unified palette)
% -----------------------------
\definecolor{BootcampBlueDark}{HTML}{1A3C68}
\definecolor{BootcampBlue}{RGB}{31,110,178}
\definecolor{AccentGold}{HTML}{DAA520}
\definecolor{SoftGray}{RGB}{245,246,248}
\definecolor{TextBlack}{HTML}{000000}
\definecolor{Crimson}{HTML}{990000}
\definecolor{MatrixGreen}{HTML}{228B22}
\definecolor{MatrixRed}{HTML}{B22222}

% -----------------------------
% Global formatting
% -----------------------------
\setbeamercolor{normal text}{fg=TextBlack,bg=white}
\setbeamercolor{title}{fg=TextBlack}
\setbeamercolor{structure}{fg=BootcampBlueDark}
\setbeamercolor{progress bar}{fg=BootcampBlueDark,bg=gray!30}

% -----------------------------
% Thick frametitle band (your “Bootcamp” look)
% -----------------------------
\setbeamercolor{frametitle}{bg=BootcampBlueDark,fg=white}
\setbeamerfont{frametitle}{series=\bfseries,size=\large}
\setbeamertemplate{frametitle}{
  \nointerlineskip
  \begin{beamercolorbox}[wd=\paperwidth,ht=3.2ex,dp=1.4ex,leftskip=1em,rightskip=1em]{frametitle}
    \usebeamerfont{frametitle}\insertframetitle
  \end{beamercolorbox}
}

% -----------------------------
% Blocks (coherent with prior lectures)
% -----------------------------
\setbeamercolor{block title example}{bg=BootcampBlueDark!90!black,fg=white}
\setbeamercolor{block body example}{bg=BootcampBlueDark!5!white,fg=black}
\setbeamertemplate{blocks}[rounded][shadow=false]

% Custom Definition Block (same API as before)
\newenvironment{definitionblock}[1]{%
  \par\vspace{0.4em}%
  \begin{beamercolorbox}[
      wd=\textwidth,
      sep=0.4em,
      leftskip=0.3em,
      rightskip=0.6em,
      rounded=false,
      shadow=false
    ]{block title example}%
    \raisebox{0.15em}{\usebeamerfont{block title example}#1}%
  \end{beamercolorbox}%
  \vspace{-0.4em}%
  \begin{beamercolorbox}[
      wd=\textwidth,
      sep=1em,
      leftskip=0.6em,
      rightskip=0.6em,
      rounded=false,
      shadow=false
    ]{block body example}%
    \usebeamerfont{block body example}%
}{%
  \end{beamercolorbox}%
  \par\vspace{0.6em}%
}

\newenvironment{proofblock}[1][Proof]{%
  \begin{block}{#1}%
}{%
  \end{block}%
}

\newcommand{\E}{\mathbb{E}}
\newcommand{\pt}{\par\noindent$\bullet$\hspace{0.6em}}
% -----------------------------
% Section / subsection transition pages
% -----------------------------
\metroset{
  sectionpage=progressbar,
  subsectionpage=progressbar
}
\AtBeginSection{
  \frame[plain]{\sectionpage}
}
\AtBeginSubsection{
  \frame[plain]{\subsectionpage}
}

% --- Theorem-like blocks (beamer)
\setbeamertemplate{theorems}[numbered] % or [ams style]

\newenvironment{theoremblock}[1]{%
  \begin{theorem}[#1]%
}{%
  \end{theorem}%
}


% Graphics folder (extracted from the provided PDF)
\graphicspath{{bm_figs/}}

% =====================================================
% Metadata
% =====================================================
\title{Option Pricing : Brownian Motion and Quadratic Variation}
\author{}
\date{}

\AtBeginSection[]{
  \begin{frame}{\insertsectionhead}
    \tableofcontents[
      currentsection,
      hideothersubsections
    ]
  \end{frame}
}

\begin{document}


% =====================================================
\begin{frame}
  \titlepage
\end{frame}

% --------------------------------------
\begin{frame}{References (chapters/pages)}

\footnotesize
\textbf{Main references.}


\begin{itemize}\setlength{\itemsep}{0.35em}
  \item \textit{Shreve}, \emph{Stochastic Calculus for Finance II}:
  Ch.~3 (Brownian Motion), pp.~83--117.
  \item \textit{Karatzas \& Shreve}, \emph{Brownian Motion}:
  Ch~2, pp.~47--126.
  \item \textit{Tomas Bj\"ork}, \emph{Arbitrage Theory in Continuous Time}:
  Ch.~3, pp.~XX--YY.
\end{itemize}

\end{frame}
% =====================================================
\begin{frame}{Roadmap}
  \tableofcontents[hideallsubsections]
\end{frame}

% ======================================================
\section{Introduction}

% ======================================================
\subsection{Examples}

% --------------------------------------
\begin{frame}{Why do we need stochastic models in finance?}

\begin{itemize}
  \item Many economic quantities evolve over time and are \textbf{uncertain}: prices, rates, demand, production, \dots \pause
  \item Decisions are taken \textbf{before} uncertainty is resolved (hedging / planning / investment). \pause
  \item We need models that capture two key empirical features: \pause
  \begin{itemize}
    \item \textbf{Randomness}: outcomes cannot be predicted deterministically. \pause
    \item \textbf{Roughness}: paths look irregular at small time scales (no classical derivative). \pause
  \end{itemize}
\end{itemize}

\medskip
\begin{alertblock}{Goal of the course}
Model randomness + roughness \;\;\(\Rightarrow\)\;\; build pricing/hedging tools for derivatives.
\end{alertblock}

\end{frame}

% --------------------------------------
% Your market examples (kept as you provided)
\begin{frame}
\frametitle{Energy prices}
{\small

\pt Airlines sell tickets several months ahead of flights, but the cost of the flight depends heavily on the kerosene price!

\vspace{5mm}

\pt Retail electricity prices for households is fixed, but market prices vary randomly depending on climate and geopolitical risks.
}

\vspace{-5mm}
\begin{figure}[!ht]
\begin{align*}
   \includegraphics[width=4.5cm]{Ressources Globales/Stochastic calculus and option pricing/Option Pricing Nizar/Lectures/Lectures 1-2/Oil.png}
   ~~~~~~~~~~~~
   \includegraphics[width=3.5cm]{Ressources Globales/Stochastic calculus and option pricing/Option Pricing Nizar/Lectures/Lectures 1-2/electricity.png}
\end{align*}
\caption{Trajectories of energy prices}
\end{figure}
\end{frame}

\begin{frame}
\frametitle{Agricultural risks}
{\small
\pt Production costs are paid ahead of time, produced quantities are revealed in the future (climate/environmental risks), and revenue depends also on market price revealed later.
}
\begin{align*}
   \includegraphics[width=5cm]{Ressources Globales/Stochastic calculus and option pricing/Option Pricing Nizar/Lectures/Lectures 1-2/ble.jpg}
   ~~~~~~~~
   \includegraphics[width=5cm]{Ressources Globales/Stochastic calculus and option pricing/Option Pricing Nizar/Lectures/Lectures 1-2/butter.jpeg}
\end{align*}
\end{frame}

\begin{frame}
\frametitle{Foreign exchange rates (FOREX)}

{\small
\pt How to face production costs in local currency with sales in foreign currency...
}

\begin{figure}[!ht]
\begin{align*}
   \includegraphics[width=5cm]{Ressources Globales/Stochastic calculus and option pricing/Option Pricing Nizar/Lectures/Lectures 1-2/EuroDollar.png}
   ~~
   \includegraphics[width=4cm]{Ressources Globales/Stochastic calculus and option pricing/Option Pricing Nizar/Lectures/Lectures 1-2/YuanDollar.jpg}
\end{align*}
\caption{Trajectories of foreign exchange rates}
\end{figure}

\end{frame}

\begin{frame}
\frametitle{Interest rates risk (fixed income)}
{\small

\pt How to manage interest rate risk?

\begin{figure}[!ht]
\begin{align*}
   \includegraphics[width=4cm]{Ressources Globales/Stochastic calculus and option pricing/Option Pricing Nizar/Lectures/Lectures 1-2/taux10ans.jpeg}
   ~~
   \includegraphics[width=4cm]{Ressources Globales/Stochastic calculus and option pricing/Option Pricing Nizar/Lectures/Lectures 1-2/TauxHypothecaires.jpg}
\end{align*}
\caption{Interest rate trajectories}
\end{figure}
}
\end{frame}

\begin{frame}
\frametitle{The stock market}
{\small
\vspace{-10mm}
\begin{figure}[!ht]
\begin{align*}
   \includegraphics[width=4cm]{Ressources Globales/Stochastic calculus and option pricing/Option Pricing Nizar/Lectures/Lectures 1-2/Apple.png}
   ~~~~~~
   \includegraphics[width=4cm]{Ressources Globales/Stochastic calculus and option pricing/Option Pricing Nizar/Lectures/Lectures 1-2/NasdaqNew.png}
\end{align*}
\caption{Trajectories of stock prices}
\end{figure}

\begin{alertblock}{}
\begin{center}{\bf
All of these curves have a very erratic nature. In this course:

how to model such erratic behavior and handle it.
}
\end{center}
\end{alertblock}
}
\end{frame}

% --------------------------------------
\begin{frame}{What kind of “random curve” are we aiming for?}

\begin{itemize}
  \item In discrete time, randomness is modeled by \textbf{coin tosses / shocks}. \pause
  \item In continuous time, we want a limit object that keeps randomness \emph{and} becomes a curve in time. \pause
  \item The canonical building block is \textbf{Brownian motion}: \pause
  \begin{itemize}
    \item continuous paths, \pause
    \item independent Gaussian increments, \pause
    \item intrinsically \textbf{rough} (not differentiable). \pause
  \end{itemize}
\end{itemize}

\medskip
\textbf{Practical takeaway:}\pause\;
in finance, we often model prices by \textbf{exponentials of Brownian motion} (log-normality) (Geometric Brownian Motion).

\end{frame}

% --------------------------------------
\begin{frame}{(Placeholder) Brownian motion pictures}
\centering

\includegraphics[width=0.47\linewidth]{Ressources Globales/Stochastic calculus and option pricing/Option Pricing Nizar/Lectures/Lectures 1-2/BM.JPG}

\medskip
\begin{itemize}
  \item Irregular at every zoom level (roughness). 
\end{itemize}
\end{frame}

% ======================================================
\subsection{Toy model: Discrete-Time model}

% --------------------------------------
\begin{frame}{A discrete-time toy world}

Fix a maturity \(T>0\) and a time grid:
\[
0=t_0<t_1<\cdots<t_N=T,
\qquad \Delta t := t_{k+1}-t_k.
\]
\pause

\begin{itemize}
  \item At each date \(t_k\), new information arrives. \pause
  \item We model information by a filtration \((\mathcal{F}_{t_k})_{k=0,\dots,N}\). \pause
  \item A price process is a sequence \((S_{t_k})_{k=0,\dots,N}\) of random variables. \pause
\end{itemize}

\medskip
\begin{block}{Key constraint (no look-ahead)}
Any trading decision at time \(t_k\) must be \(\mathcal{F}_{t_k}\)-measurable.
\end{block}

\end{frame}

% --------------------------------------
\begin{frame}{A binomial stock model (one-step)}

At each step, the stock either goes \textbf{up} or \textbf{down}:
\[
S_{t_{k+1}} =
\begin{cases}
u\,S_{t_k} & \text{(up) with probability $p$}\\
d\,S_{t_k} & \text{(down) with probability $1-p$}
\end{cases}
\qquad \text{with } 0<d<1<u.
\]
\pause

\begin{itemize}
  \item This is the cleanest setting to understand replication and hedging. \pause
  \item The risk-free asset grows deterministically:
  \[
  B_{t_{k+1}}=(1+r\Delta t)\,B_{t_k}.
  \]
  \pause
\end{itemize}

\medskip
\textbf{Next:}\pause\;
draw the tree for \(N=2,3\), then price claims by backward induction.

\end{frame}

% --------------------------------------
\begin{frame}{Tree picture for \(N=2\) (template)}

\centering
\begin{tikzpicture}[scale=0.95]
\node (S0) at (0,0) {\(S_0\)};
\node (Su) at (2,1.2) {\(uS_0\)};
\node (Sd) at (2,-1.2) {\(dS_0\)};
\node (Suu) at (4,2.2) {\(u^2S_0\)};
\node (Sud) at (4,0) {\(udS_0\)};
\node (Sdd) at (4,-2.2) {\(d^2S_0\)};

\draw[->] (S0) -- (Su) node[midway, above] {\small up};
\draw[->] (S0) -- (Sd) node[midway, below] {\small down};
\draw[->] (Su) -- (Suu) node[midway, above] {\small up};
\draw[->] (Su) -- (Sud) node[midway, below] {\small down};
\draw[->] (Sd) -- (Sud) node[midway, above] {\small up};
\draw[->] (Sd) -- (Sdd) node[midway, below] {\small down};
\end{tikzpicture}

\medskip
\begin{itemize}
  \item A claim payoff is a function of the terminal node(s). \pause
  \item Pricing will become a \textbf{backward recursion}. 
\end{itemize}

\end{frame}

% --------------------------------------
\begin{frame}{From many tosses to a continuous-time limit (idea)}

To approach continuous time, we accelerate the steps:\pause
\begin{itemize}
  \item \(n\) steps per unit time, so \([0,t]\) contains \(nt\) tosses, \pause
  \item step size becomes of order \(1/\sqrt{n}\), so the CLT applies. \pause
\end{itemize}

\medskip
A canonical scaling for a symmetric random walk \((X_k)\) is:\pause
\[
B^{(n)}(t) := \frac{1}{\sqrt{n}}\,S_{\lfloor nt\rfloor}
\quad
\text{(with linear interpolation when needed).}
\]
\pause

\begin{block}{Message}
As \(n\to\infty\), the scaled random walk behaves like a Brownian motion.
\end{block}

\end{frame}

% --------------------------------------
\begin{frame}{From binomial prices to geometric Brownian motion}

A standard “diffusion scaling” for the binomial model is:\pause
\[
u_n = 1+\frac{\sigma}{\sqrt{n}},
\qquad
d_n = 1-\frac{\sigma}{\sqrt{n}},
\qquad
\text{(with \(n\) steps per unit time).}
\]
\pause

The resulting stock price at time \(t\) can be written (heads/tails count):\pause
\[
S_n(t) = S(0)\,u_n^{H_{nt}}\,d_n^{T_{nt}}.
\]
\pause

\begin{alertblock}{Limit (heuristic)}
As \(n\to\infty\), \(S_n(t)\) converges (in distribution) to
\[
S(t)=S(0)\exp\!\Big(\sigma W_t - \tfrac12\sigma^2 t\Big),
\]
i.e. a \textbf{geometric Brownian motion}.
\end{alertblock}

\end{frame}

% --------------------------------------
\begin{frame}{Takeaways of the introduction}

\begin{itemize}
  \item Real financial curves are \textbf{random} and \textbf{rough}. \pause
  \item Continuous-time limits lead to \textbf{Brownian motion} and \textbf{geometric Brownian motion}.
\end{itemize}

\end{frame}

% =====================================================
% =====================================================
% INTRODUCTION (completed)
% =====================================================
% ======================================================
\section{Brownian Motion}

% ======================================================
\subsection{Random Walk}

% --------------------------------------
\begin{frame}{Random Walk (Discrete Time)}

\begin{definitionblock}{Simple Symmetric Random Walk}
Let \((\varepsilon_k)_{k\ge 1}\) be i.i.d. with
\[
\mathbb{P}(\varepsilon_k=+1)=\mathbb{P}(\varepsilon_k=-1)=\tfrac12.
\]\pause
Define
\[
S_0=0,
\qquad
S_n := \sum_{k=1}^n \varepsilon_k,\quad n\ge 1.
\]
\end{definitionblock}\pause

\begin{itemize}
  \item \((S_n)\) is a discrete-time stochastic process.\pause
  \item Increments are \(\Delta S_n = S_n-S_{n-1}=\varepsilon_n\).\pause
  \item Natural filtration: \(\mathcal{F}_n:=\sigma(\varepsilon_1,\dots,\varepsilon_n)=\sigma(S_0,\dots,S_n)\).
\end{itemize}

\end{frame}

% --------------------------------------
\begin{frame}{Illustrations (Random Walk Paths)}

\begin{itemize}
  \item Different outcomes \(\omega\) produce different sample paths \(n\mapsto S_n(\omega)\).\pause
  \item Typical behavior:\pause
  \begin{itemize}
    \item oscillates up and down,\pause
    \item wanders away from 0 at scale \(\sqrt{n}\).\pause
  \end{itemize}
\end{itemize}

\medskip
\begin{center}
% (Replace with your figure, or keep as a placeholder.)
\includegraphics[width=0.45\linewidth]{Ressources Globales/Stochastic calculus and option pricing/Option Pricing Nizar/Lectures/Lectures 1-2/random-walk.png}
\end{center}

\end{frame}

% --------------------------------------
\begin{frame}{Basic Properties and Continuous Interpolation}

\begin{itemize}
  \item \textbf{Mean and variance:}\pause
  \[
  \mathbb{E}[S_n]=0,
  \qquad
  \mathrm{Var}(S_n)=n.
  \]
  \item \textbf{Independent increments:}\pause
  \(S_{n+m}-S_n = \sum_{k=n+1}^{n+m}\varepsilon_k\) is independent of \(\mathcal{F}_n\).\pause
  \item \textbf{Martingale property:}\pause
  \[
  \mathbb{E}[S_{n+1}\mid \mathcal{F}_n]=S_n.
  \]
\end{itemize}
\pause
\medskip

\textbf{Continuous interpolation:}\pause\;
define a continuous-time process \(\widetilde S(t)\) by linear interpolation:
\[
\widetilde S(t) := S_k + (t-k)(S_{k+1}-S_k),
\quad t\in[k,k+1].
\]

\end{frame}

% --------------------------------------
\begin{frame}{Rescaled Continuous Interpolation}

We now introduce a small time step \(1/n\) and accelerate time.\pause

\medskip
\textbf{Rescaled process:}\pause
\[
W^{(n)}(t)
:=
\frac{1}{\sqrt{n}}\,
\widetilde S(nt),
\qquad t\in[0,T].
\]
\pause

\begin{itemize}
  \item Time is sped up by a factor \(n\) (many steps per unit time).\pause
  \item Space is scaled by \(1/\sqrt{n}\) (CLT scaling).\pause
  \item Typical fluctuations over \([0,t]\) stay of order \(\sqrt{t}\).\pause
\end{itemize}

\medskip
\begin{block}{Goal}
Understand the limit of \(W^{(n)}\) as \(n\to\infty\).
\end{block}

\end{frame}

% --------------------------------------
\begin{frame}{Behavior of Increments (Heuristic)}

Fix \(0\le s<t\).\pause

\medskip
The increment of the rescaled walk is approximately:
\[
W^{(n)}(t)-W^{(n)}(s)
\approx
\frac{1}{\sqrt{n}}\Big(S_{\lfloor nt\rfloor}-S_{\lfloor ns\rfloor}\Big).
\]
\pause

\begin{itemize}
  \item It is a sum of \(\lfloor n(t-s)\rfloor\) i.i.d.\ \(\pm1\) variables.\pause
  \item Mean \(0\), variance \(\approx t-s\).\pause
  \item By the CLT, it should look Gaussian:\pause
  \[
  W^{(n)}(t)-W^{(n)}(s)\ \approx\ \mathcal{N}(0,t-s).
  \]
\end{itemize}

\medskip
\textbf{Key intuition:}\pause\;
many tiny independent shocks \(\Rightarrow\) Gaussian fluctuations.

\end{frame}

% --------------------------------------
\begin{frame}{Convergence as the Time Step Goes to 0 (Donsker)}

\begin{theoremblock}{Donsker's Invariance Principle (Informal)}
As \(n\to\infty\), the rescaled interpolated random walk \(W^{(n)}\) converges in distribution
to a continuous-time process \((W_t)_{t\ge 0}\) called \textbf{Brownian motion}.
\end{theoremblock}\pause

\begin{itemize}
  \item This is the continuous-time analogue of the CLT.\pause
  \item Brownian motion is the universal scaling limit of many random walks.\pause
\end{itemize}

\medskip
\textbf{Interpretation:}\pause\;
Brownian motion is the canonical model of continuous-time randomness.

\end{frame}

% ======================================================
\subsection{Definition of Brownian Motion}

% --------------------------------------
\begin{frame}{Definition of Brownian Motion}

\begin{definitionblock}{Standard Brownian Motion}
A process \((W_t)_{t\ge 0}\) is a \emph{standard Brownian motion} if:\pause
\begin{itemize}
  \item \(W_0=0\),\pause
  \item \textbf{independent increments},\pause
  \item \textbf{Gaussian increments:} for \(0\le s<t\),
  \[
  W_t-W_s \sim \mathcal{N}(0,t-s),\pause
  \]
  \item \textbf{continuous paths} \(t\mapsto W_t(\omega)\) (a.s.).
\end{itemize}
\end{definitionblock}

\end{frame}



% --------------------------------------
\begin{frame}{Gaussian Density and the Heat Equation}

For each \(t>0\), the random variable \(W_t\) admits a density:
\[
W_t \sim \mathcal{N}(0,t)
\qquad\Longrightarrow\qquad
p(t,x)
=
\frac{1}{\sqrt{2\pi t}}
\exp\!\Big(-\frac{x^2}{2t}\Big).
\]
\pause

\medskip
\begin{itemize}
  \item More generally, for \(0\le s<t\),
  \[
  W_t \mid W_s=x \sim \mathcal{N}(x,\ t-s),
  \quad\text{with transition density}\quad
  p(t-s,y-x).
  \]\pause
  \item The density \(p(t,x)\) solves the \textbf{heat equation}:
  \[
  \partial_t p(t,x)
  =
  \frac12\,\partial_{xx}p(t,x),
  \qquad t>0,\ x\in\mathbb{R},
  \]
  with initial condition \(p(0,x)=\delta_0(x)\) (Dirac mass at \(0\)).\pause
\end{itemize}

\medskip
\begin{block}{Preview (PDE link)}
Many pricing PDEs (e.g.\ Black--Scholes) are perturbations of the heat equation.
\end{block}

\end{frame}


% --------------------------------------
\begin{frame}{Brownian Motion as a Gaussian Process}

\begin{itemize}
  \item Brownian motion is a \textbf{Gaussian process}:\pause
  any finite vector \((W_{t_1},\dots,W_{t_m})\) is multivariate normal.\pause
  \item It is fully characterized by:\pause
  \begin{itemize}
    \item mean function \(m(t)=\mathbb{E}[W_t]=0\),\pause
    \item covariance kernel \(K(s,t)=\mathrm{Cov}(W_s,W_t)=\min(s,t).\pause\)
  \end{itemize}
\end{itemize}

\medskip
\begin{block}{Key consequence}
For \(0\le s<t\),
\[
W_t = W_s + (W_t-W_s),
\]
where \(W_s\) and \(W_t-W_s\) are independent Gaussian variables.
\end{block}

\end{frame}

% --------------------------------------
\begin{frame}{Transformations of Brownian Motion}

Let \((W_t)_{t\ge0}\) be a Brownian motion.\pause

\begin{itemize}
  \item \textbf{Symmetry:}\pause\;
  \((-W_t)_{t\ge0}\) is also Brownian motion.\pause
  \item \textbf{Scaling:}\pause\;
  for any \(c>0\),
  \[
  \widetilde W_t := \frac{1}{\sqrt{c}}W_{ct}
  \quad \text{is Brownian motion.}
  \]
  \pause
  \item \textbf{Time inversion (advanced but useful):}\pause\;
  the process
  \[
  \widehat W_t := t\, W_{1/t}, \qquad t>0,
  \]
  is again a Brownian motion (in law).\pause
  \item \textbf{Time reversal on \([0,T]\):}\pause\;
  \[
  \big(W_T - W_{T-t}\big)_{0\le t\le T}
  \quad \text{is Brownian motion.}
  \]
\end{itemize}

\end{frame}

% --------------------------------------
\begin{frame}{Continuity (What to Assume vs What to Prove?)}

\begin{itemize}
  \item One can \textbf{define} Brownian motion by finite-dimensional distributions (Gaussian increments).\pause
  \item A deep theorem (Kolmogorov) guarantees existence of a \textbf{continuous version}.\pause
  \item For this course, we will:\pause
  \begin{itemize}
    \item take continuity as part of the definition,\pause
    \item and focus on consequences for calculus and finance.\pause
  \end{itemize}
\end{itemize}

\medskip
\begin{block}{Practical viewpoint}
The main “new” phenomenon is not continuity itself,
but the \textbf{roughness} of paths (non-differentiability, quadratic variation).
\end{block}

\end{frame}

% ======================================================
\subsection{Markovian Properties of Brownian Motion}

\subsubsection{Weak Markov property }
% --------------------------------------
\begin{frame}{Markov Property (What it means)}

\begin{definitionblock}{Markov Property (informal)}
A process is Markov if \textbf{the future depends on the past only through the present}.
\end{definitionblock}\pause

\medskip
For Brownian motion, for \(s<t\):\pause
\[
\mathbb{P}(W_t\in A \mid \mathcal{F}_s)
=
\mathbb{P}(W_t\in A \mid W_s),
\]
for all Borel sets \(A\).\pause

\medskip
\begin{itemize}
  \item The increment \(W_t-W_s\) is independent of \(\mathcal{F}_s\).\pause
  \item The conditional law is Gaussian:\pause
  \[
  W_t \mid W_s \sim \mathcal{N}(W_s,\ t-s).
  \]
\end{itemize}

\end{frame}

% --------------------------------------
\begin{frame}{Weak Markov Property}

Let $t \geq 0$.

\begin{definitionblock}{Weak Markov Property}
Conditionally on \(\mathcal{F}_t\), the shifted process
\[
\big(W_{s+t}-W_t\big)_{s\ge0}
\]
is a Brownian motion independent of \(\mathcal{F}_t\).
\end{definitionblock}\pause

\begin{itemize}
  \item The Markov property holds even if the “present time” is random.\pause
  \item This is fundamental for hitting times, barrier events, and American options.
\end{itemize}

\end{frame}

% --------------------------------------
\begin{frame}{Weak Markov Property (Test-Functions Formulation)}

% Fix \(t\ge 0\).\pause

\begin{definitionblock}{Weak Markov Property (via test functions)}
For any \(n\in\mathbb{N}\), any times \(0\le s_1<\cdots<s_n\), and any bounded measurable
\(\varphi:\mathbb{R}^n\to\mathbb{R}\), we have almost surely
\[
\mathbb{E}\!\left[
\varphi\!\big(W_{t+s_1}-W_t,\dots,W_{t+s_n}-W_t\big)
\,\Big|\,\mathcal{F}_t
\right]
=
\mathbb{E}\!\left[
\varphi\!\big(W_{s_1},\dots,W_{s_n}\big)
\right].
\]
\end{definitionblock}\pause

\begin{itemize}
  \item Future increments after time \(t\) have the \emph{same law} as a fresh Brownian motion.\pause
  \item The conditional law given \(\mathcal{F}_t\) is independent of the past.\pause
  \item Strong Markov: the same statement with \(t\) replaced by a stopping time \(\tau\).
\end{itemize}

\end{frame}

% --------------------------------------
\begin{frame}{Weak Markov Property}

\begin{definitionblock}{Weak Markov Property}
For \(0\le s<t\) and any bounded measurable \(f\),\pause
\[
\mathbb{E}\big[f(W_t)\mid \mathcal{F}_s\big]
=
u_{t-s}(W_s),
\]
where
\[
u_{h}(x)
:=
\mathbb{E}\big[f(x+Z)\big],
\quad Z\sim\mathcal{N}(0,h).
\]
This is equivalent to say that $W_t - W_s$
\end{definitionblock}\pause

\begin{itemize}
  \item The conditional expectation depends on the past only through \(W_s\).\pause
  \item This is already enough to derive many transition formulas and PDE links.
\end{itemize}

\end{frame}

% --------------------------------------
\begin{frame}{Brownian Motion is a Martingale}

Let $(W_t)_{t\ge 0}$ be a standard Brownian motion and let
\[
\mathcal{F}_t:=\sigma(W_s:\ 0\le s\le t)
\qquad\text{(natural filtration).}
\]

\begin{theoremblock}{Martingale property}
$(W_t)_{t\ge 0}$ is a martingale with respect to $(\mathcal{F}_t)_{t\ge 0}$, i.e.
for all $0\le s\le t$,
\[
\mathbb{E}\!\left[\,W_t \mid \mathcal{F}_s\,\right]=W_s
\qquad\text{a.s.}
\]
\end{theoremblock}\pause

\begin{proofblock}{Proof (one line)}
Write $W_t=W_s+(W_t-W_s)$; the increment $(W_t-W_s)$ is independent of $\mathcal{F}_s$
and centered, hence
\[\mathbb{E}\!\left[W_t\mid\mathcal{F}_s\right]
=
W_s+\mathbb{E}\!\left[W_t-W_s\mid\mathcal{F}_s\right]
=
W_s.
\]
\end{proofblock}
\end{frame}

\subsubsection{Strong Markov Property}
% --------------------------------------
\begin{frame}{Strong Markov Property (Stopping Times)}

Let \(\tau\) be a stopping time w.r.t.\ \((\mathcal{F}_t)\).\pause

\begin{definitionblock}{Strong Markov Property}
Conditionally on \(\mathcal{F}_\tau\), the shifted process
\[
\big(W_{\tau+t}-W_\tau\big)_{t\ge0}
\]
is a Brownian motion independent of \(\mathcal{F}_\tau\).
\end{definitionblock}\pause

\begin{itemize}
  \item The Markov property holds even if the “present time” is random.\pause
  \item This is fundamental for hitting times, barrier events, and American options.
\end{itemize}

\end{frame}


% --------------------------------------
\begin{frame}{Brownian Motion is a Martingale}

Let $(W_t)_{t\ge 0}$ be a standard Brownian motion and let
\[
\mathcal{F}_t:=\sigma(W_s:\ 0\le s\le t)
\qquad\text{(natural filtration).}
\]

\begin{theoremblock}{Martingale property}
$(W_t)_{t\ge 0}$ is a martingale with respect to $(\mathcal{F}_t)_{t\ge 0}$, i.e.
for all $0\le s\le t$,
\[
\mathbb{E}\!\left[\,W_t \mid \mathcal{F}_s\,\right]=W_s
\qquad\text{a.s.}
\]
\end{theoremblock}\pause

\begin{proofblock}{Proof (one line)}
Write $W_t=W_s+(W_t-W_s)$; the increment $(W_t-W_s)$ is independent of $\mathcal{F}_s$
and centered, hence
\[\mathbb{E}\!\left[W_t\mid\mathcal{F}_s\right]
=
W_s+\mathbb{E}\!\left[W_t-W_s\mid\mathcal{F}_s\right]
=
W_s.
\]
\end{proofblock}
\end{frame}
% --------------------------------------
\begin{frame}{Strong Markov Property (Test-Functions Formulation)}

 Let \(\tau\) be a stopping time w.r.t.\ \((\mathcal{F}_t)\).\pause

\begin{definitionblock}{Strong Markov Property (via test functions)}
 For any \(n\in\mathbb{N}\), any times \(0\le t_1<\cdots<t_n\), and any bounded measurable
\(\varphi:\mathbb{R}^n\to\mathbb{R}\), we have almost surely
\[
\mathbb{E}\!\left[
\varphi\!\big(W_{\tau+t_1}-W_\tau,\dots,W_{\tau+t_n}-W_\tau\big)
\,\Big|\,\mathcal{F}_\tau
\right]
=
\mathbb{E}\!\left[
\varphi\!\big(W_{t_1},\dots,W_{t_n}\big)
\right].
\]
\end{definitionblock}\pause

% \begin{itemize}
%   \item The future increments after \(\tau\) have the \emph{same law} as a fresh Brownian motion.\pause
%   \item The conditional law given \(\mathcal{F}_\tau\) does \emph{not} depend on the past.\pause
%   \item Key tool for hitting times, barrier events, and American options.\pause
% \end{itemize}

\end{frame}



% ======================================================
% ======================================================
\subsection{Reflection Principle}

% --------------------------------------
\begin{frame}{Reflection Principle: Why it matters}

\begin{itemize}
  \item Define the running maximum
  \[
  M_t := \sup_{0\le s\le t} W_s .
  \]
  \item The pair \((M_t,W_t)\) is central for \textbf{barrier options} and for \textbf{default / first-passage} modeling. 
  \item Key idea: events of the form \(\{M_t\ge a\}\) are the same as hitting events
  \[
  \{M_t\ge a\} \iff \{\tau_a \le t\},
  \qquad
  \tau_a := \inf\{s\ge 0:\ W_s=a\}. 
  \]
\end{itemize}

\medskip
\begin{block}{What we gain}
The reflection principle converts \(\{M_t\ge a,\,W_t\le w\}\) into a \emph{simple Gaussian tail event},
hence explicit formulas.
\end{block}

\end{frame}



% --------------------------------------
\begin{frame}{Reflection Map (Informal Geometric Picture)}

Fix \(a>0\) and \(t>0\). Consider a path \(\omega\) such that it hits level \(a\) before \(t\)
(i.e. \(\tau_a(\omega)\le t\)). Define a reflected path \(\omega^\star\) by
\[
\omega^\star(s)=
\begin{cases}
\omega(s), & 0\le s\le \tau_a(\omega),\\
2a-\omega(s), & \tau_a(\omega) < s\le t.
\end{cases}
\]
\medskip

\begin{itemize}
  \item \(\omega^\star\) coincides with \(\omega\) up to the first hit of \(a\), then mirrors it around \(a\).
  \item Reflection swaps “ending below \(a\)” with “ending above \(a\)” among paths that \emph{did hit} \(a\).
\end{itemize}


\end{frame}


% --------------------------------------
\begin{frame}{Reflection Map (Why Wiener measure is preserved)}

Fix \(a>0\) and \(t>0\). For a continuous path \(\omega:[0,t]\to\mathbb{R}\), define
\[
\tau_a(\omega):=\inf\{s\in[0,t]:\ \omega(s)=a\}
\quad (\inf\emptyset:=+\infty).
\]
On the event \(\{\tau_a(\omega)\le t\}\), define the reflected path \(\omega^\star\) by
\[
\omega^\star(s)=
\begin{cases}
\omega(s), & 0\le s\le \tau_a(\omega),\\
2a-\omega(s), & \tau_a(\omega) < s\le t.
\end{cases}
\]

\medskip
\begin{itemize}
  \item \textbf{Geometric effect:}\pause\;
  after the first hit of level \(a\), we mirror the future trajectory around the horizontal line \(y=a\).\pause

  \item \textbf{Involution:}\pause\;
  applying the map twice gives back the original path:
  \((\omega^\star)^\star=\omega\).

  \end{itemize}

\end{frame}


\begin{frame}{Why Wiener measure is preserved (informal but key): }
    
  \begin{itemize}
    \item Brownian increments are \textbf{symmetric}: for any \(h>0\),
    \(\Delta W := W_{u+h}-W_u\) has the same law as \(-\Delta W\).\pause
    \item Conditioned on the past up to \(\tau_a\), the future increments
    \((W_{\tau_a+s}-W_{\tau_a})_{0\le s\le t-\tau_a}\)
    have the same law as their sign-flipped version.\pause
    \item Reflection is exactly this sign flip after \(\tau_a\), shifted to keep continuity at level \(a\):
    \(W_{\tau_a+s}\mapsto 2a-W_{\tau_a+s}\).\pause
    \item Hence the mapping \(\omega\mapsto\omega^\star\) sends Brownian paths to Brownian paths
    with the \textbf{same probability weight} (Wiener measure is invariant under this reflection).\pause
\item \textbf{Consequence:}\pause\;
  on the event \(\{\tau_a\le t\}\), reflection sends \(W_t\) to \(W_t^\star=2a-W_t\), hence it swaps
  \[
  \{W_t\le a\}\quad \longleftrightarrow \quad \{W_t^\star\ge a\},
  \]
  (up to the null boundary \(\{W_t=a\}\)), which is the key step behind the reflection principle.
\end{itemize}

\end{frame}

% --------------------------------------
\begin{frame}{Reflection Principle (Event-Level Statement)}

Let \(a>0\), \(t>0\), and \(w\le a\). Then the reflection principle yields
\[
\mathbb{P}\big(M_t > a,\ W_t < w\big)
=
\mathbb{P}\big(W_t > 2a-w\big),
\qquad (w\le a). \]

\pause
\medskip
\begin{block}{Most useful special case}
Taking \(w=a\) gives
\[
\mathbb{P}(M_t \ge a)
=
2\,\mathbb{P}(W_t \ge a).
\]
\end{block}

\end{frame}

% --------------------------------------
\begin{frame}{Joint Law of \texorpdfstring{$(M_t,W_t)$}{(M_t,W_t)}}

A powerful refinement is the explicit joint density of \((M_t,W_t)\):  )
\[
f_{M_t,W_t}(m,w)
=
\frac{2(2m-w)}{t\sqrt{2\pi t}}
\exp\!\Big(-\frac{(2m-w)^2}{2t}\Big),
\qquad (m>0,\ w\le m). 
\]

\pause
\medskip
\begin{itemize}
  \item It comes from differentiating the reflection identity
  \[
  \mathbb{P}(M_t\ge m,\ W_t\le w)=\mathbb{P}(W_t\ge 2m-w),
  \quad w\le m.  
  \]
  \item From this, you can derive many useful quantities:
  \(\mathcal L(M_t)\), conditional laws \(M_t\mid W_t=w\), barrier probabilities, etc.  
\end{itemize}

\end{frame}

% --------------------------------------
\begin{frame}{Finance Link: Barrier Events = Maxima / Hitting Times}

\begin{itemize}
  \item A (toy) up-and-out barrier event typically looks like
  \[
  \{\text{option survives to }t\}
  =
  \{M_t < H\}.
  \]
  \item Credit/default via first-passage often models default as
  \[
  \tau_H := \inf\{t\ge 0:\ X_t \le H\}
  \quad \text{or} \quad
  \inf\{t\ge 0:\ X_t \ge H\},
  \]
  i.e. again a barrier / hitting time.
\end{itemize}

\end{frame}

\subsection{Pathwise Properties of Brownian Motion}

% --------------------------------------
\begin{frame}{Oscillations at Every Time}

Let \((W_t)_{t\ge 0}\) be Brownian motion.\pause

\begin{theoremblock}{Local oscillation at scale \(\sqrt{h}\)}
For every fixed \(t\ge 0\), with probability \(1\),
\[
\limsup_{h\downarrow 0}\frac{W_{t+h}-W_t}{\sqrt{h}}=+\infty,
\qquad
\liminf_{h\downarrow 0}\frac{W_{t+h}-W_t}{\sqrt{h}}=-\infty.
\]
\end{theoremblock}\pause

\begin{itemize}
  \item Brownian motion crosses above and below its current level infinitely often on any interval.\pause
  \item No flat pieces, no monotone pieces, and no local trend that persists at small scales.\pause
\end{itemize}

\medskip
\begin{theoremblock}{Nowhere differentiability}
With probability \(1\), the path \(t\mapsto W_t(\omega)\) is nowhere differentiable.
\end{theoremblock}

\end{frame}

% --------------------------------------
\begin{frame}{Behavior at Infinity: \texorpdfstring{$\limsup=+\infty$}{limsup=+infty} and \texorpdfstring{$\liminf=-\infty$}{liminf=-infty}}

\begin{theoremblock}{Recurrence and unbounded oscillations {\color{blue}\textbf{(See Homework)}}}
With probability \(1\),
\[
\limsup_{t\to\infty} W_t = +\infty,
\qquad
\liminf_{t\to\infty} W_t = -\infty.
\]
and 
\[
\lim_{t\to\infty} \frac{W_t}{t} = 0,
\]
\end{theoremblock}\pause

\begin{itemize}
  \item Brownian motion does not drift to \(+\infty\) or \(-\infty\): it keeps coming back and overshooting.\pause
  \item In fact, it hits every level infinitely many times (recurrence in 1D).\pause
\end{itemize}

\medskip
\begin{block}{Interpretation}
Over long horizons, fluctuations grow like \(\sqrt{t}\), so the process eventually reaches arbitrarily large positive and negative values.
\end{block}

\end{frame}

% --------------------------------------
\begin{frame}{Law of the Iterated Logarithm (Sharp Envelope)}

\begin{theoremblock}{LIL (as \(t\to\infty\))}
With probability \(1\),
\[
\limsup_{t\to\infty}\frac{W_t}{\sqrt{2t\log\log t}} = 1,
\qquad
\liminf_{t\to\infty}\frac{W_t}{\sqrt{2t\log\log t}} = -1.
\]
\end{theoremblock}\pause

\begin{itemize}
  \item \(\sqrt{t}\) gives the correct scale, but the \(\sqrt{\log\log t}\) factor is the \emph{sharp} correction.\pause
  \item This is the precise “envelope” of Brownian motion at infinity.
\end{itemize}


\end{frame}

% --------------------------------------
\begin{frame}{Modulus of Continuity (Small-Time Envelope)}

\begin{theoremblock}{LIL near 0 (modulus of continuity)}
If $\alpha<1/2$, then almost surely Brownian motion
is everywhere locally $\alpha$ Holder continuous : 
\[
  |W_t(\omega)-W_s(\omega)|
  \le
  C_{\alpha,T}(\omega)\,|t-s|^{\alpha},
  \qquad \forall\,s,t\in[0,T].
  \]
  \pause
 For \(\alpha=\tfrac12\), this fails: Brownian motion is \emph{not} Holder-\(\tfrac12\).\pause
\end{theoremblock}

\medskip
\begin{block}{Interpretation}
Brownian motion is ``almost'' \(1/2\)-Holder: it is Holder for every exponent strictly below \(1/2\),
but never as regular as \(1/2\) itself.
\end{block}

\end{frame}

% --------------------------------------
\begin{frame}{Hitting Times and Level Crossings (Pathwise Consequence)}

For \(a\in\mathbb{R}\), define the hitting time
\[
\tau_a := \inf\{t\ge 0:\ W_t=a\}.
\]\pause

\begin{theoremblock}{Almost sure hitting (1D) {\color{blue}\textbf{(See Homework)}} }
For any \(a\in\mathbb{R}\),\quad \(\mathbb{P}(\tau_a<\infty)=1\).
\end{theoremblock}\pause

\begin{itemize}
  \item Brownian motion hits every level with probability 1.\pause
  \item Combined with independent increments, it implies repeated crossings and strong oscillations.\pause
  \item This is the starting point for barrier options and first-passage computations.
\end{itemize}

\end{frame}


% ======================================================
\section{Quadratic Variation}

% ======================================================
\subsection{Motivation}

% --------------------------------------
\begin{frame}{Why Quadratic Variation? (Financial Motivation)}

\textbf{Goal:} define the gain of a \emph{self-financing} trading strategy in continuous time.\pause

\medskip
In discrete time (rebalancing dates \(0=t_0<t_1<\dots<t_n=T\)):\pause
\[
\text{Gain}(T)
=
\sum_{k=0}^{n-1} \Delta_k \big(S_{t_{k+1}}-S_{t_k}\big),
\]
where \(\Delta_k\) is the number of shares held on \([t_k,t_{k+1})\).\pause

\medskip
\textbf{Continuous-time limit:}\pause\;
we want to make sense of
\[
\int_0^T \Delta_t \, dS_t
\qquad
\text{as a limit of discrete sums.}
\]

\medskip
\begin{alertblock}{Key difficulty}
If \(S\) is driven by Brownian motion, its paths are \emph{very rough} (nowhere differentiable).\pause\\
Usual Riemann--Stieltjes integrals do \emph{not} apply.
\end{alertblock}

\end{frame}

% % --------------------------------------
% \begin{frame}{A Toy Warning: ``\texorpdfstring{$dW_t$}{dW} is of order \texorpdfstring{$\sqrt{dt}$}{sqrt(dt)}''}

% Let \(W\) be Brownian motion.\pause
% Over a small time step \(\Delta t\):\pause
% \[
% W_{t+\Delta t}-W_t \sim \mathcal{N}(0,\Delta t)
% \quad\Rightarrow\quad
% |W_{t+\Delta t}-W_t|\ \text{typically of size } \sqrt{\Delta t}.
% \]\pause

% \medskip
% So if we try to treat \(W\) as differentiable,\pause
% \[
% \frac{W_{t+\Delta t}-W_t}{\Delta t}
% \sim
% \frac{\sqrt{\Delta t}}{\Delta t}
% =
% \frac{1}{\sqrt{\Delta t}}
% \to \infty,
% \]
% which hints that ``\(\dot W_t\)'' does not exist.\pause

% \medskip
% \textbf{But:} squared increments behave nicely!\pause
% \[
% (W_{t+\Delta t}-W_t)^2
% \ \text{is of size } \Delta t.
% \]

% \medskip
% This is the origin of \textbf{quadratic variation}.

% \end{frame}

% --------------------------------------
\begin{frame}{From Discrete Trading to Stochastic Integrals}

\textbf{Self-financing idea (informal).}\pause\;
A portfolio \((V_t)\) holding \(\Delta_t\) shares and the rest in cash satisfies:\pause
\[
dV_t \approx \Delta_t\, dS_t
\quad\text{(no external cash injection).}
\]\pause

\medskip
To make \(\int_0^T \Delta_t\, dS_t\) rigorous when \(S\) is Brownian-driven,\pause
we need a limit theory for sums of the form\pause
\[
\sum_{k} \Delta_{t_k}\big(S_{t_{k+1}}-S_{t_k}\big).
\]\pause

\medskip
\begin{itemize}
  \item The limit will be the \textbf{It\^o integral} (defined first for simple strategies).\pause
  \item Its second-moment behavior is governed by \textbf{quadratic variation}.\pause
\end{itemize}

\medskip
\textbf{Punchline:}\pause\;
quadratic variation is the mechanism behind the extra ``\(dt\)'' term in It\^o calculus.

\end{frame}

% ======================================================
\subsection{Definition of Quadratic Variation}

% --------------------------------------
\begin{frame}{Quadratic Variation: Definition Along a Partition}

Let \(X=(X_t)_{0\le t\le T}\) be a real-valued process.\pause

\medskip
A partition \(\pi\) of \([0,T]\) is\pause
\[
\pi:\ 0=t_0<t_1<\dots<t_n=T,
\qquad
|\pi|:=\max_k (t_{k+1}-t_k).
\]\pause

\begin{definitionblock}{Quadratic Variation Along \(\pi\)}
Define the quadratic sum
\[
[X]^{\pi}_T
:=
\sum_{k=0}^{n-1}\big(X_{t_{k+1}}-X_{t_k}\big)^2.
\]
\end{definitionblock}\pause

If \([X]^{\pi}_T\) converges (in probability, or in \(L^2\)) as \(|\pi|\to 0\),
we call the limit the \textbf{quadratic variation} of \(X\) on \([0,T]\), denoted \([X]_T\).

\end{frame}

% --------------------------------------
\begin{frame}{First Examples: Smooth vs.\ Rough}

\textbf{Smooth case (finite variation).}\pause\;
If \(X_t=\int_0^t v_s\,ds\) with \(v\) integrable,\pause
then increments are of order \(\Delta t\), hence\pause
\[
[X]^{\pi}_T
=
\sum_k O((\Delta t_k)^2)
\ \longrightarrow\ 0.
\]\pause

\medskip
So\ \(\boxed{[X]_T=0}\) for finite-variation processes.\pause

\medskip
\textbf{Rough case (random walk intuition).}\pause\;
For a nearest-neighbor random walk \(S_i=\sum_{j=1}^i Y_j\) with \(Y_j\in\{\pm1\}\),\pause
\[
\sum_{j=1}^i (S_j-S_{j-1})^2
=
\sum_{j=1}^i Y_j^2
=
i,
\]
a deterministic growth.\pause

\medskip
Brownian motion is the continuous-time analogue: squared increments accumulate to time.

\end{frame}

% --------------------------------------
\begin{frame}{Quadratic Variation of Brownian Motion: The Statement}

Let \(W\) be a standard Brownian motion.\pause

\begin{theoremblock}{Quadratic Variation of \(W\)}
For any \(T>0\), as \(|\pi|\to 0\),
\[
[W]^{\pi}_T
=
\sum_{k=0}^{n-1}\big(W_{t_{k+1}}-W_{t_k}\big)^2
\ \longrightarrow\ T
\quad
\text{in } L^2\ \text{(hence in probability).}
\]
In other words,\quad \(\boxed{[W]_T=T}\).
\end{theoremblock}\pause

\medskip
\textbf{Interpretation:}\pause
\begin{itemize}
  \item Brownian paths are not differentiable.\pause
  \item Yet their squared increments ``average out'' to elapsed time.
\end{itemize}

\end{frame}

% --------------------------------------
\begin{frame}{A Key Consequence: A Martingale Built from \texorpdfstring{$W^2$}{W2}}

Quadratic variation explains why \(W_t^2\) is \emph{not} a martingale.\pause

\medskip
Compute for \(s<t\):\pause
\[
\mathbb{E}\big[W_t^2 \mid \mathcal{F}_s\big]
=
\mathbb{E}\big[(W_s+(W_t-W_s))^2 \mid \mathcal{F}_s\big]
=
W_s^2 + \mathbb{E}\big[(W_t-W_s)^2\big].
\]\pause

But \(\mathbb{E}[(W_t-W_s)^2]=t-s\), so\pause
\[
\mathbb{E}\big[W_t^2 \mid \mathcal{F}_s\big]
=
W_s^2 + (t-s).
\]\pause

\begin{block}{Conclusion}
\[
\boxed{\,W_t^2 - t \text{ is a martingale.}\,}
\]
This is the simplest glimpse of It\^o’s formula.
\end{block}

\end{frame}



% --------------------------------------
\begin{frame}{Quadratic (Co)Variation: A Useful Extension}

For two processes \(X,Y\), define along \(\pi\):\pause
\[
[X,Y]^{\pi}_T
:=
\sum_{k=0}^{n-1}\big(X_{t_{k+1}}-X_{t_k}\big)\big(Y_{t_{k+1}}-Y_{t_k}\big).
\]\pause

\medskip
\textbf{Examples.}\pause
\begin{itemize}
  \item If \(Y=X\), we recover \([X]^{\pi}_T\).\pause
  \item If \(A\) has finite variation and \(W\) is Brownian, then typically\pause
  \[
  [A,W]_T=0.
  \]
\end{itemize}

\medskip
\begin{block}{Finance preview}
In multi-factor models, \([W^{(i)},W^{(j)}]_t\) encodes instantaneous correlation,
and drives hedging and risk aggregation.
\end{block}

\end{frame}

% --------------------------------------
\begin{frame}{Preview: Why Quadratic Variation Creates It\^o Terms}

Heuristic rule for Brownian increments:\pause
\[
(\Delta W)^2 \approx \Delta t,
\qquad
\Delta W \cdot \Delta t \approx 0,
\qquad
(\Delta t)^2 \approx 0.
\]\pause

\medskip
So when expanding \(f(W_{t+\Delta t})\) by Taylor,\pause
\[
f(W+\Delta W)
\approx
f(W) + f'(W)\Delta W + \tfrac12 f''(W)(\Delta W)^2,
\]\pause
and since \((\Delta W)^2\) accumulates to \(dt\), we expect\pause
\[
df(W_t)
=
f'(W_t)\,dW_t
+
\frac12 f''(W_t)\,dt.
\]\pause

\begin{alertblock}{Key message}
Quadratic variation is the rigorous reason why stochastic calculus
is \emph{not} ordinary calculus.
\end{alertblock}

\end{frame}

% --------------------------------------
\begin{frame}{Heuristic Notation: $dW_t$ and $dt$}

Think of $\Delta W_k$ as a small increment $dW$ over a small time step $\Delta t_k=dt$.\pause

\medskip
\begin{block}{Quadratic variation as a rule of computation}
The limit $Q_t^\pi=\sum (\Delta W)^2 \to t$ is encoded as
\[
(dW_t)^2 = dt.
\]
\end{block}\pause

\begin{block}{Mixed and time terms are negligible}
The limits
\[
\sum \Delta W\,\Delta t \to 0,
\qquad
\sum (\Delta t)^2 \to 0
\]
are encoded as
\[
dW_t\,dt = 0,
\qquad
(dt)^2 = 0.
\]
\end{block}\pause

\begin{itemize}
  \item This is \textbf{not} algebra on infinitesimals: it is a compact way to remember
  which discrete sums vanish and which ones accumulate.\pause
  \item These rules will become rigorous once the It\^o integral and quadratic variation
  are defined.
\end{itemize}

\end{frame}

% --------------------------------------
\begin{frame}{The It\^o Multiplication Table (One-Dimensional)}

Whenever you expand expressions and keep terms up to order $dt$, use:\pause

\begin{theoremblock}{It\^o table (formal but reliable)}
\[
\boxed{
(dW_t)^2 = dt,
\qquad
dW_t\,dt = 0,
\qquad
(dt)^2 = 0.
}
\]
\end{theoremblock}\pause

\medskip
\begin{itemize}
  \item Rule of thumb: $\Delta W = O(\sqrt{\Delta t})$, so
  $(\Delta W)^2 = O(\Delta t)$ but $\Delta W\,\Delta t = O((\Delta t)^{3/2})$.\pause
  \item Hence, in Taylor expansions, \textbf{keep} $(\Delta W)^2$ but \textbf{drop}
  $\Delta W\,\Delta t$ and $(\Delta t)^2$.
\end{itemize}

\end{frame}


% ======================================================
\section{Geometric Brownian Motion (GBM)}

% ======================================================
\subsection{Motivation: why multiplicative models?}

% --------------------------------------
\begin{frame}{Why Not Model Prices with an Additive Brownian Motion?}

A first naive idea would be
\[
S_t = S_0 + \mu t + \sigma W_t.
\]\pause

\begin{itemize}
  \item \textbf{Problem 1 (negativity):}\pause\;
  an additive Gaussian model can make \(S_t<0\), which is unrealistic for most assets.\pause
  \item \textbf{Problem 2 (scale):}\pause\;
  a \$1 move does not mean the same thing for a \$10 stock vs a \$1000 stock.\pause
  \item \textbf{Empirical cue:}\pause\;
  returns are closer to being \emph{scale-free} than prices.
\end{itemize}

\medskip
\begin{block}{Idea}
Model \textbf{relative changes} (returns) rather than absolute changes.
\end{block}

\end{frame}

% --------------------------------------
\begin{frame}{Log-Returns: Additive and (Approximately) Gaussian}

Define the \textbf{log-price} and \textbf{log-return}:
\[
X_t := \log S_t,
\qquad
\log\frac{S_{t+\Delta t}}{S_t} = X_{t+\Delta t}-X_t.
\]\pause

\begin{itemize}
  \item Log-returns are additive over time: \(\log(S_T/S_0)=\sum \log(S_{t_{k+1}}/S_{t_k})\).\pause
  \item In simple models, log-returns have \textbf{stationary increments}:\pause\;
  distribution depends mainly on \(\Delta t\).\pause
  \item This naturally suggests:
  \[
  S_t = S_0 e^{ (\mu - \frac{\sigma^2}{2})t + \sigma W_t}
  \]
\end{itemize}

\end{frame}

\subsection{Distributional properties}

% --------------------------------------
\begin{frame}{Lognormal Distribution (Marginal and Conditional Laws)}

Since \(W_t\sim\mathcal N(0,t)\), we obtain\pause
\[
\log S_t \sim \mathcal N\!\Big(
\log S_0 + (\mu-\tfrac12\sigma^2)t,\ \sigma^2 t
\Big).
\]\pause

Hence \(S_t\) is \textbf{lognormal}.\pause

\medskip
\textbf{Conditional distribution (Markov + independent increments):}\pause\;
for \(0\le s<t\),
\[
\log\frac{S_t}{S_s}
=
\Big(\mu-\tfrac12\sigma^2\Big)(t-s)
+
\sigma (W_t-W_s),
\]
so\pause
\[
\log\frac{S_t}{S_s}\ \Big|\ \mathcal F_s
\sim
\mathcal N\!\Big(
(\mu-\tfrac12\sigma^2)(t-s),\ \sigma^2(t-s)
\Big).
\]

\end{frame}

% --------------------------------------
\begin{frame}{Moments and Volatility Scaling (Quick Facts)}

From the explicit form,\pause
\[
\mathbb E[S_t]
=
S_0 e^{\mu t},
\qquad
\mathrm{Var}(S_t)
=
S_0^2 e^{2\mu t}\big(e^{\sigma^2 t}-1\big).
\]\pause

\begin{itemize}
  \item Over a short interval \(\Delta t\), log-return satisfies\pause
  \[
  \log\frac{S_{t+\Delta t}}{S_t}
  \approx
  \mathcal N\!\Big((\mu-\tfrac12\sigma^2)\Delta t,\ \sigma^2\Delta t\Big).
  \]\pause
  \item This is the origin of the common rule: \(\text{vol} \propto \sqrt{\Delta t}\).
\end{itemize}

\end{frame}

% ======================================================
\subsection{Calibration (Geometric Brownian Motion)}

% --------------------------------------
\begin{frame}{Calibration of GBM: from data to \texorpdfstring{$\sigma$}{sigma} (historical volatility)}

Assume the GBM model
\[
S(t)=S(0)\exp\!\Big((\mu-\tfrac12\sigma^2)t+\sigma W_t\Big),
\]
so that \textbf{log-returns} satisfy, for a partition \(0=t_0<\cdots<t_n=T\),\pause
\[
R_j
:=\log\!\frac{S_{t_{j+1}}}{S_{t_j}}
=
\sigma\big(W_{t_{j+1}}-W_{t_j}\big)
+
\Big(\mu-\tfrac12\sigma^2\Big)(t_{j+1}-t_j).
\]

\medskip
\begin{block}{Idea: use quadratic variation}
Since \([W]_T=T\), the leading term in \(\sum_j R_j^2\) is
\[
\sum_{j=0}^{n-1} R_j^2
\;\approx\;
\sigma^2\sum_{j=0}^{n-1}(W_{t_{j+1}}-W_{t_j})^2
\;\approx\;
\sigma^2 T,
\]
while the cross-terms and drift terms vanish as \(|\pi|\to 0\).\pause
\end{block}

\begin{alertblock}{Practical estimator (realized variance)}
\[
\widehat{\sigma}^2
=
\frac{1}{T}\sum_{j=0}^{n-1}
\Big(\log\!\frac{S_{t_{j+1}}}{S_{t_j}}\Big)^2,
\qquad
\widehat{\sigma}=\sqrt{\widehat{\sigma}^2}.
\]
(If data are daily, take \(t\) in \emph{years} to annualize automatically.)
\end{alertblock}

\end{frame}

% --------------------------------------
\begin{frame}{Calibration: historical vs implied (risk-neutral)}

\begin{itemize}
  \item \textbf{Historical calibration (under \( \mathbb{P}\))}: estimate \(\sigma\) (and sometimes \(\mu\)) from a time series of prices using log-returns and realized variance.\pause
  \item \textbf{Implied calibration (under \( \mathbb{Q}\))}: choose \(\sigma\) so that Black--Scholes prices match observed option prices (implied volatility surface).\pause
\end{itemize}

\medskip
\begin{block}{Important message for pricing}
\begin{itemize}
  \item In no-arbitrage pricing, \(\mu\) disappears and is replaced by \(r\) under \(\mathbb{Q}\).\pause
  \item In practice, \(\sigma\) from \(\mathbb{P}\) (historical) and \(\sigma\) from \(\mathbb{Q}\) (implied) typically differ.\pause
\end{itemize}
\end{block}

\medskip
\textbf{Takeaway:}\pause\;
\(\sigma\) is the key parameter to \emph{calibrate}.

\end{frame}

% ======================================================



% --------------------------------------
\begin{frame}{What GBM Gets Right / Wrong}

\begin{itemize}
  \item \textbf{Gets right (as a first model):}\pause
  positivity, tractability, volatility scaling, closed forms.\pause
  \item \textbf{Gets wrong (empirically):}\pause
  constant volatility, Gaussian tails, no jumps, no volatility clustering.\pause
  \item Still, GBM is the \textbf{baseline} on top of which richer models are built.
\end{itemize}


\end{frame}


\section*{Discussion on stopping times and stopped Martingales}


% --------------------------------------
\begin{frame}{Stopping Times and Stopped Processes}

Let $(M_t)_{t\ge 0}$ be an $(\mathcal F_t)$-martingale and let $\tau$ be a stopping time.\pause

\begin{definitionblock}{Stopped process}
For $t\ge 0$, define
\[
M^{\tau}_t := M_{t\wedge \tau}.
\]
\end{definitionblock}\pause

\begin{itemize}
  \item $M^{\tau}$ follows $M$ up to time $\tau$, then stays constant.\pause
  \item If $\tau$ is finite a.s., then $M^\tau_t$ converges a.s. to $M_\tau$ as $t\to\infty$.\pause
\end{itemize}

\medskip
\begin{block}{Goal}
Show that stopping a martingale at a stopping time preserves the martingale property.
\end{block}

\end{frame}

% --------------------------------------
\begin{frame}{Optional Stopping: A Basic Martingale Preservation Result}

Let $(M_t)_{t\ge 0}$ be an $(\mathcal F_t)$-martingale.\pause

\begin{theoremblock}{Stopped martingale}
For any stopping time $\tau$, the stopped process $(M_{t\wedge\tau})_{t\ge 0}$
is an $(\mathcal F_t)$-martingale, i.e.\ for all $0\le s\le t$,
\[
\E\!\left[M_{t\wedge\tau}\mid \mathcal F_s\right]=M_{s\wedge\tau}
\qquad\text{a.s.}
\]
\end{theoremblock}\pause

\begin{proofblock}{Proof (clean decomposition)}
Decompose the event $\{\tau\le s\}$ vs.\ $\{\tau>s\}$:\pause
\[
M_{t\wedge\tau}
=
M_\tau\,\mathbf 1_{\{\tau\le s\}}
+
M_{t\wedge\tau}\,\mathbf 1_{\{\tau>s\}}.
\]
On $\{\tau\le s\}$, we have $t\wedge\tau=s\wedge\tau=\tau$, hence the conditional expectation is $M_\tau$.

\end{proofblock}

\end{frame}

% --------------------------------------
\begin{frame}{Optional Stopping Theorem (OST): When Can We Replace $t\wedge\tau$ by $\tau$?}

The stopped-process result holds for \emph{any} stopping time.\pause

\medskip
But to conclude something about $\E[M_\tau]$ (or $\E[M_\tau\mid \mathcal F_s]$),
we need integrability/uniform integrability assumptions.\pause

\begin{theoremblock}{Optional Stopping (typical safe versions)}
Let $(M_t)$ be a martingale and $\tau$ a stopping time.\pause

\begin{itemize}
  \item If $\tau$ is \textbf{bounded} (i.e.\ $\tau\le T$ a.s.), then
  \[
  \E[M_\tau]=\E[M_0].
  \]\pause

  \item If $\tau$ is \textbf{integrable} and $M$ has \textbf{bounded increments} (discrete time),
  then still $\E[M_\tau]=\E[M_0]$.\pause

  \item If $(M_{t\wedge\tau})_{t\ge 0}$ is \textbf{uniformly integrable} (UI), then
  \[
  \E[M_\tau]=\E[M_0]
  \quad\text{and}\quad
  \E[M_\tau\mid \mathcal F_s]=M_{s\wedge\tau}.
  \]
\end{itemize}
\end{theoremblock}

\begin{alertblock}{Warning}
``$\tau<\infty$ a.s.'' alone is \textbf{not} enough to justify $\E[M_\tau]=\E[M_0]$.
\end{alertblock}

\end{frame}

% --------------------------------------
\begin{frame}{Why the Warning Matters (A Brownian Motion Example)}

Let $(W_t)$ be Brownian motion, a martingale.\pause

\medskip
Take the hitting time $\tau_a:=\inf\{t\ge 0:\ W_t=a\}$ for $a>0$.\pause

\begin{itemize}
  \item We know $\mathbb{P}(\tau_a<\infty)=1$.\pause
  \item But $\E[\tau_a]=+\infty$ (heavy tail), and $W_{\tau_a}=a$ a.s.\pause
\end{itemize}

\medskip
If we (incorrectly) applied $\E[W_{\tau_a}]=\E[W_0]$, we would get $a=0$, absurd.\pause

\begin{block}{Correct statement}
For each $t$, optional stopping is safe for the bounded time $t\wedge\tau_a$:
\[
\E[W_{t\wedge\tau_a}] = \E[W_0]=0.
\]
But taking $t\to\infty$ requires UI / integrability assumptions.
\end{block}

\end{frame}

% --------------------------------------
\begin{frame}{A Handy result: $M_{t\wedge\tau}$ is a Martingale in Practice}

Let $(M_t)$ be a martingale and $\tau$ a stopping time.\pause

\begin{itemize}
  \item For any $t\ge 0$,
  \[
  \E[M_{t\wedge\tau}] = \E[M_0].
  \]
  (Just take $s=0$ in the stopped-martingale property.)\pause

  \item If $M$ is nonnegative (or bounded below), then $(M_{t\wedge\tau})_{t\ge 0}$
  is a supermartingale-dominated family, and one can often pass to limits using Fatou.\pause
\end{itemize}

\end{frame}


% --------------------------------------
\begin{frame}{Laplace Transform of the First Hitting Time}

Let $(W_t)_{t\ge0}$ be a Brownian motion and, for $a\in\mathbb{R}$,
$
\tau_a := \inf\{t\ge0:\ W_t=a\}.
$
\vspace{-6mm}
\pause 
\begin{theoremblock}{Laplace transform of $\tau_a$}
For every $\lambda>0$,
\vspace{-3mm}
\[
\E\!\left[e^{-\lambda \tau_a}\right]
=
\exp\!\left(-|a|\sqrt{2\lambda}\right).
\vspace{-6mm}
\]
\end{theoremblock}
\pause

\textbf{Sketch of computation (exponential martingale).}\pause
For $\theta\in\mathbb{R}$, define
\[
M_t^{(\theta)} := \exp\!\Big(\theta W_t-\tfrac12\theta^2 t\Big),
\qquad t\ge0,
\vspace{-3mm}
\]
then $(M_t^{(\theta)})_{t\ge0}$ is a martingale. \pause
Apply optional stopping to $\tau_a\wedge t$:
\[
1=\E[M_{\tau_a\wedge t}^{(\theta)}]
=
\E\!\left[\exp\!\Big(\theta W_{\tau_a\wedge t}-\tfrac12\theta^2(\tau_a\wedge t)\Big)\right].
\vspace{-3mm}
\]
\pause
On $\{\tau_a\le t\}$, $W_{\tau_a}=a$, hence
\[
1 \ge \E\!\left[\exp\!\Big(\theta a-\tfrac12\theta^2\tau_a\Big)\mathbf{1}_{\{\tau_a\le t\}}\right].
\vspace{-3mm}
\]
Let $t\to\infty$ (monotone convergence) and choose $\theta=\sqrt{2\lambda}$ to get
\[
\E\!\left[e^{-\lambda \tau_a}\right]\le e^{-\sqrt{2\lambda}\,a}\quad(a>0).
\vspace{-3mm}
\]
\pause
Apply the same argument to $-W$ to recover the $|a|$ and obtain the stated formula.

\end{frame}

\end{document}
